\documentclass[UTF8]{ctexart}
\usepackage[a4paper,margin=2.5cm]{geometry}
\usepackage{booktabs}
\usepackage{graphicx}
\usepackage{hyperref}
\graphicspath{{report/}{./}}
\title{基于轻量化网络与知识蒸馏的 3D 点云目标检测方法设计与分析}
\author{人工智能课程大作业}
\date{2026年6月}
\begin{document}
\begin{center}\LARGE 基于轻量化网络与知识蒸馏的 3D 点云目标检测方法设计与分析\end{center}

\section{摘要}

3D 点云目标检测是自动驾驶、机器人感知和三维场景理解中的重要任务。与普通二维图像相比，点云数据具有稀疏、不规则和空间结构明显等特点，因此检测模型通常需要较复杂的特征提取结构。复杂模型可以提升检测效果，但也会带来参数量大、计算量高、部署困难等问题。本文以 Li 等人在 Applied Sciences 2023 发表的 A Lightweight Model for 3D Point Cloud Object Detection 为参考[1]，借鉴其中轻量化稀疏卷积和知识蒸馏的思想，设计了一个面向课程实验的 3D 体素目标检测方案。本文没有把目标设定为照搬原论文工程系统，而是在课程条件下围绕“基线实现、轻量化改进、蒸馏增强、结果比较”建立完整实验流程。


本文首先构建了一个小规模合成 3D 体素数据集，将目标分为 car、pedestrian 和 cyclist 三类，并为每个样本设置类别标签和归一化三维框。基线模型采用标准特征提取结构，用于提供基本检测能力；创新实验一将大规模特征变换替换为分组变换和通道混合结构，以减少模型复杂度；创新实验二在轻量模型基础上加入知识蒸馏，使学生模型学习基线模型的输出分布、中间特征和特征关系。实验结果显示，轻量模型参数量约压缩到基线的 1/14.2，单样本乘加量约压缩到 1/14.3；加入知识蒸馏后，分类准确率由 83.75\% 提升到 89.58\%，BBox MAE 由 0.1337 降低到 0.1241。结果说明，在模型规模显著降低的情况下，蒸馏训练可以帮助轻量模型恢复部分性能。


关键词：3D 点云；目标检测；轻量化模型；知识蒸馏；课程实验


\section{1 绪论}

\subsection{1.1 研究背景}

点云由激光雷达等传感器采集得到，能够描述物体在三维空间中的位置和几何结构。在自动驾驶场景中，车辆需要识别周围的汽车、行人和骑行者，并给出它们的三维框位置。如果检测模型太大，即使精度较高，也很难部署到车载或边缘设备上。因此，如何在尽量保持检测效果的同时降低参数量和计算量，是一个比较实际的问题。


在三维检测任务中，模型不仅要判断目标类别，还要预测目标在三维空间中的位置和尺度。与二维目标检测相比，三维框多了深度、高度和空间方向等信息，数据表达也更复杂。点云通常是稀疏的，很多空间位置没有点；同时点的数量和分布会受传感器距离、遮挡、反射材料等因素影响。因此，三维检测模型往往需要较强的特征提取能力。


但是，强特征提取能力经常意味着更大的模型。对于课程实验来说，如果直接选择一个完整的自动驾驶检测框架，不仅数据量很大，而且还会涉及体素化、anchor 生成、稀疏卷积库、后处理和官方评测协议等工程细节。本文更关注一个清晰的问题：在一个可控的小规模 3D 检测任务中，能否先实现一个基线模型，再用轻量化结构降低复杂度，并用知识蒸馏弥补轻量化带来的性能下降。


\subsection{1.2 本文工作内容}

本文围绕参考论文中的核心思想开展课程实验，但重点是设计自己的方案并进行验证。具体来说，本文首先设计了一个小规模 3D 体素目标检测任务，使模型同时完成类别预测和三维框回归；然后实现标准基线模型，作为后续比较的基础；在此基础上设计轻量化学生模型，用分组变换和通道混合减少参数量与计算量；随后在轻量模型上加入知识蒸馏，包括输出蒸馏、特征蒸馏和关系蒸馏；最后从定量指标和定性效果图两个角度比较基线实验与创新实验。


\subsection{1.3 参考论文的作用}

本文参考的论文提出了 LW-Sconv 模块和三段式知识蒸馏[1]。论文中的基线是 SECOND[2]，提升方法是在 3D sparse convolution 中加入分组卷积、因子化卷积和深度卷积思想，再让轻量学生模型学习教师模型的输出、特征和特征流关系[1]。该论文对本文的主要启发不是某一个具体代码细节，而是一个研究思路：先建立可靠基线，再从模型结构和训练策略两个方向降低复杂度并恢复性能。


\section{2 研究现状}

三维目标检测方法大致可以分为点方法、体素方法和混合方法。PointNet 直接处理点集合，VoxelNet 和 SECOND 将点云划分成体素后用卷积提取特征，PointPillars 则把点云划分成柱状结构后投影到鸟瞰图。随着模型越来越深，检测效果通常会提高，但计算量也明显增加。论文中提到的 SECOND 已经使用稀疏卷积提升效率，不过仍然有较高 FLOPs[2]。


轻量化模型设计在二维图像任务中比较常见，例如 MobileNet 使用 depth-wise separable convolution[4]，ShuffleNet 使用 group convolution 和 channel shuffle。论文的主要贡献就是把这些轻量化思想迁移到 3D sparse convolution 中[1]。同时，知识蒸馏让小模型学习大模型的“软知识”，不仅学习真实标签，也学习教师模型的输出分布和中间特征[3]。


\subsection{2.1 基于点的检测方法}

基于点的方法直接以点集作为输入，典型代表是 PointNet 及其后续方法。这类方法的优点是保留了点云原始结构，不需要把点云强行转换成规则网格。但点云是无序集合，点的数量也不固定，因此模型需要设计对点排列不敏感的特征聚合方式。PointNet 使用多层感知机提取每个点的特征，再通过最大池化得到全局特征。这种方法结构直观，但对局部几何关系的刻画能力有限。


\subsection{2.2 基于体素的检测方法}

基于体素的方法先把三维空间划分为规则网格，再把点云映射到体素中。VoxelNet 和 SECOND 属于这一类。体素化之后，模型可以使用卷积结构提取空间特征，整体流程更接近传统卷积神经网络。缺点是三维体素网格可能非常稀疏，如果使用普通三维卷积，会在大量空位置上浪费计算。SECOND 通过稀疏卷积减少无效计算，是三维检测中比较有代表性的基线方法。


\subsection{2.3 基于 BEV 表示的方法}

还有一些方法将点云投影到鸟瞰图，也就是 BEV 表示。PointPillars 将点云划分为柱状结构，再投影到二维平面上处理。BEV 表示的优点是计算更接近二维卷积，速度较快，也便于自动驾驶场景中理解车辆周围平面布局。缺点是高度信息会被压缩，需要通过特征编码尽量保留三维结构。


\subsection{2.4 模型轻量化与知识蒸馏}

模型轻量化的目标是在尽量少损失性能的情况下减少参数量、计算量和存储开销。常见方式包括剪枝、量化、低秩分解、深度可分离卷积、分组卷积和通道混合等。本文主要借鉴分组变换和通道混合思想[1]。知识蒸馏则是一种训练策略，它让小模型学习大模型的输出分布和中间特征[3]。与只学习硬标签相比，蒸馏提供的信息更丰富。例如教师模型可能认为某个样本 80\% 像汽车、15\% 像骑行者、5\% 像行人，这种软分布比单纯的 one-hot 标签更能体现类别之间的相似关系。


\section{3 方法}

\subsection{3.1 任务定义与总体思路}

本文将实验任务定义为一个小规模 3D 体素目标检测问题。输入是固定大小的三维体素网格，输出包括目标类别和归一化三维框参数。目标类别包括 car、pedestrian 和 cyclist，三维框使用 center\_z、center\_y、center\_x、size\_z、size\_y、size\_x 表示。这个任务虽然比真实自动驾驶点云检测简单，但仍保留了类别识别和空间定位两个核心目标。


参考论文给本文的主要启发是：三维检测模型不一定只能通过堆叠更复杂的网络提升效果，也可以通过轻量化结构降低复杂度，再用知识蒸馏恢复部分性能[1]。因此，本文方法分为两条线：第一条是基线实验，用标准模型验证任务可学习性；第二条是创新实验，在轻量模型上加入蒸馏训练，与基线进行比较。


\subsection{3.2 基线实验}

（a）输入与标签。基线实验的输入为体素化后的 3D 网格。每个样本包含一个目标，标签由目标类别和三维框参数组成。由于体素网格大小固定，三维框参数被归一化到 0 到 1，便于模型学习。


（b）模型结构。基线模型采用标准特征提取结构：首先将三维体素网格转换为特征向量，然后经过标准特征提取层得到中间表示，最后分别接分类头和框回归头。分类头输出三类目标的预测结果，框回归头输出 6 个三维框参数。基线模型参数量较大，表达能力较强，适合作为对照模型。


（c）训练方式。基线模型采用监督学习训练，只使用真实标签进行约束。分类部分使用交叉熵损失，框回归部分使用均方误差。总损失可以写为 Lsup = Lcls + lambda\_box Lbox，其中 lambda\_box 用来控制框回归损失的权重。基线实验流程如图1所示。


\begin{figure}[htbp]\centering
\includegraphics[width=0.95\linewidth]{assets/baseline\_flow.png}
\caption{基线实验流程图}
\end{figure}

\subsection{3.3 创新实验}

（a）轻量化学生模型。创新实验首先将基线模型中的大规模特征变换替换为分组变换和通道混合结构。分组变换相当于让不同特征组分别进行计算，可以减少参数量和乘加量；通道混合则用于缓解分组后信息不流通的问题，使不同组之间的特征重新融合。


（b）知识蒸馏增强。轻量模型虽然复杂度低，但容量变小后容易损失性能。因此，本文进一步加入知识蒸馏。蒸馏时，基线模型作为教师模型，轻量模型作为学生模型[3]。学生模型不仅学习真实标签，也学习教师模型的输出分布、中间特征和特征关系[1]。这样可以让轻量模型在参数较少的情况下获得更多训练信息。


（c）创新实验与基线的区别。基线实验只使用标准模型和真实标签；创新实验则同时改变模型结构和训练方式。结构上，创新实验使用轻量化分组特征提取；训练上，创新实验加入教师模型提供的软标签和中间特征监督。创新实验总体流程如图2所示。


\begin{figure}[htbp]\centering
\includegraphics[width=0.95\linewidth]{assets/method\_flow.png}
\caption{创新实验总体方法流程图}
\end{figure}

\subsection{3.4 损失函数与评价指标}

本文模型同时进行分类和三维框回归。分类损失用于约束目标类别预测，框回归损失用于约束目标位置和大小。知识蒸馏部分包括三项：Lout 约束学生模型输出接近教师模型，Lfeature 约束学生模型中间特征接近教师模型，Lrelation 约束学生模型样本间特征关系接近教师模型。综合来看，创新实验的目标函数可以表示为：L = Lsup + alpha\_out Lout + alpha\_feature Lfeature + alpha\_relation Lrelation。


评价指标包括参数量、单样本乘加量、分类准确率、BBox MAE 和 mean IoU。参数量和乘加量衡量模型复杂度；分类准确率衡量类别判断能力；BBox MAE 衡量三维框参数误差；mean IoU 衡量预测框与真实框的空间重叠程度。由于这些指标关注角度不同，本文在结果分析中同时使用定量表格、指标图和定性效果图。


\section{4 实验与结果}

\subsection{4.1 数据集划分}

本文共生成 960 个样本，其中 720 个用于训练，240 个用于验证。训练集和验证集使用不同随机种子生成，避免模型只记住某一批固定样本。每个样本包含一个目标，目标类别在三类之间随机选择。由于本文任务是课程级方法验证，数据规模没有追求很大，而是强调每个样本标签清晰、结构可控、能够稳定比较不同模型。


\subsection{4.2 模型设置}

实验设置三个模型。第一是 Baseline teacher，对应基线实现。第二是 Lightweight student，对应仅使用轻量化结构的学生模型。第三是 Lightweight + KD，对应加入知识蒸馏训练的轻量学生模型。三个模型使用相同的数据集和评价指标，保证比较具有一致性。


\subsection{4.3 评价方式}

评价分为定量和定性两部分。定量评价包括参数量、单样本乘加量、分类准确率、BBox MAE 和 mean IoU。参数量和乘加量用于衡量模型复杂度，分类准确率用于衡量类别判断能力，BBox MAE 和 mean IoU 用于衡量空间定位能力。定性评价通过 BEV 投影图展示真实框和预测框的位置关系，能够直观看出模型预测是否贴近目标。


\subsection{4.4 定量结果}

\begin{center}\small\begin{tabular}{lrrrrr}\hline
模型 & 参数量 & 单样本乘加量 & 分类准确率 & BBox MAE & mean IoU \\
\hline
Baseline teacher & 114089 & 113952 & 87.08\% & 0.1020 & 0.2012 \\
Lightweight student & 8033 & 7968 & 83.75\% & 0.1337 & 0.1191 \\
Lightweight + KD & 8033 & 7968 & 89.58\% & 0.1241 & 0.1146 \\
\end{tabular}\end{center}

\begin{figure}[htbp]\centering
\includegraphics[width=0.95\linewidth]{assets/results\_visual.png}
\caption{复杂度和性能指标可视化}
\end{figure}

从表格和图3可以看到，轻量模型的参数量从 114089 降到 8033，单样本乘加量从 113952 降到 7968，压缩幅度比较明显。按参数量计算，轻量模型约为基线的 1/14.2；按乘加量计算，轻量模型约为基线的 1/14.3。这说明分组瓶颈结构确实达到了降低复杂度的目标。


从检测效果看，普通轻量模型分类准确率为 83.75\%，低于基线模型的 87.08\%，说明单纯压缩模型会带来一定性能损失。加入知识蒸馏后，Lightweight + KD 的分类准确率提升到 89.58\%，超过普通轻量模型，也略高于基线模型。BBox MAE 方面，普通轻量模型为 0.1337，蒸馏模型降至 0.1241，说明蒸馏对框参数平均误差也有一定改善。


mean IoU 的变化没有分类准确率那么理想。蒸馏模型 mean IoU 为 0.1146，略低于普通轻量模型的 0.1191。这说明本文蒸馏策略主要改善了类别判断和框参数平均误差，但对框的空间重叠质量提升不足。原因可能是本文使用的框回归损失是参数级误差，而不是直接优化 IoU，因此模型可能在中心和尺寸的平均误差上变小，但局部样本的重叠效果没有同步提升。


\subsection{4.5 定性效果图}

为了更直观地观察检测效果，本文将验证样本投影到 BEV 平面，并同时绘制真实框和预测框。绿色框表示真实框，红色框表示模型预测框。图4选取 car、pedestrian 和 cyclist 三类代表样本，对比了基线模型、普通轻量模型和蒸馏模型的预测效果。


\begin{figure}[htbp]\centering
\includegraphics[width=0.95\linewidth]{assets/qualitative\_results.png}
\caption{定性检测效果图}
\end{figure}

从图4可以看到，模型不仅给出类别预测，也给出了对应的目标位置。对于 car 和 cyclist 这类形状较明显的样本，蒸馏模型的预测框与真实框重叠较好，说明学生模型能够从教师模型中学习到一定定位能力。对于 pedestrian 这类较小目标，预测仍然更困难，因为目标在 BEV 平面上的占用区域很小，少量误差就会导致 IoU 明显下降。这也解释了为什么本文整体 mean IoU 不如分类准确率好看。


\subsection{4.6 基线与创新实验比较}

基线模型的优势是结构直接、容量较大、训练稳定。它在定位指标上表现最好，mean IoU 达到 0.2012，说明较大的模型容量对框回归有帮助。但基线模型的缺点也很明显：参数量和计算量都远高于轻量模型。


普通轻量模型的优势是复杂度最低，参数量只有 8033，乘加量只有 7968。它的问题是单纯压缩后性能下降，分类准确率和框回归误差都不如基线。这个结果符合轻量化模型的常见规律：模型变小以后，表达能力会受到限制。


蒸馏增强模型在轻量模型基础上加入教师指导，分类准确率和 BBox MAE 都优于普通轻量模型。它说明知识蒸馏可以缓解轻量化带来的部分性能损失。不过，蒸馏模型的 mean IoU 没有超过普通轻量模型，说明当前蒸馏方式还没有完全解决空间定位质量问题。


\section{5 总结}

本文围绕轻量化 3D 点云目标检测完成了一次课程实验。报告按照课程要求，先介绍任务背景和研究现状，再给出方法设计，随后进行实验验证和结果分析。本文没有把参考论文作为照搬对象，而是借鉴其中轻量化结构和知识蒸馏的思想[1]，设计了基线实验和创新实验，并对二者进行了比较。


实验结果表明，轻量化学生模型能够显著减少参数量和单样本乘加量，但单纯压缩会带来一定性能下降；加入知识蒸馏后，轻量模型的分类准确率和 BBox MAE 得到改善，说明教师模型提供的软标签和中间特征对学生模型有帮助。定性效果图也显示，模型可以在 BEV 投影中给出较直观的目标框预测。


本文仍有不足。由于课程时间和算力限制，实验使用的是小规模合成体素数据，而不是真实 KITTI 或 nuScenes 数据集；模型结构也没有实现完整稀疏卷积检测框架。因此，本文结果主要用于验证方法思路，而不能直接代表真实自动驾驶场景效果。总体来看，本次实验完成了基线实现、方法优化和方案验证，符合课程报告中对实际内容和方法思考的要求。


\section{代码可用性声明}

本文相关实验代码已整理并开源至 GitHub，仓库地址为：https://github.com/ShaoZidi/ai-course-3d-pointcloud-lightweight-kd 。


\section{参考文献}

[1] Li, Z.; Li, Y.; Wang, Y.; Xie, G.; Qu, H.; Lyu, Z. A Lightweight Model for 3D Point Cloud Object Detection. Applied Sciences, 2023, 13(11), 6754.


[2] Yan, Y.; Mao, Y.; Li, B. SECOND: Sparsely Embedded Convolutional Detection. Sensors, 2018.


[3] Hinton, G.; Vinyals, O.; Dean, J. Distilling the Knowledge in a Neural Network. arXiv, 2015.


[4] Howard, A. G. et al. MobileNets: Efficient Convolutional Neural Networks for Mobile Vision Applications. arXiv, 2017.

\end{document}
